\chapter{Literature Survey}

This chapter provides a comprehensive review of the fundamental technologies and recent research related to multi-camera object tracking and re-identification systems designed for city-scale video surveillance. The discussion is organized into four main sections. Section 3.1 establishes the technical foundations of object detection and single-camera tracking. Section 3.2 examines how systems identify and match individuals and vehicles across different camera views using visual features. Section 3.3 reviews recent research from 2022 to 2025, comparing different approaches and identifying their strengths and limitations. Finally, Section 3.4 explains the specific technologies chosen for this project, providing clear justifications based on performance, practicality, and suitability for Egyptian Ministry of Interior operations.

\section{Background on Object Detection and Single-Camera Tracking}

Modern surveillance systems rely on two core tasks: identifying objects in each video frame and following those objects over time. This section explains the foundational approaches.

\subsection{Object Detection Using Deep Learning}

Object detection locates and classifies objects within images. In surveillance applications, the primary targets are people and vehicles. The You Only Look Once (YOLO) family represents the dominant approach in real-time systems. Unlike two-stage detectors that first propose regions then classify them, YOLO performs detection in a single forward pass, enabling real-time performance.

YOLO26m (Ultralytics) is the detection model used in this system. It achieves [PLACEHOLDER]\% mean Average Precision (mAP) on the COCO benchmark with [PLACEHOLDER] parameters. Architectural improvements in the YOLO26m design include enhanced feature extraction through an improved backbone and neck, and a spatial pyramid pooling fusion (SPPF) module that significantly boosts detection of small and occluded objects common in urban surveillance — partially obscured vehicles at intersections, overlapping pedestrians in crowds, and distant objects at the periphery of wide-angle cameras.

\subsection{Single-Camera Multi-Object Tracking}

After detecting objects in each frame, the tracking task connects detections across frames into continuous trajectories. The tracking-by-detection paradigm is standard: a detector identifies all objects in frame $t$, and the tracker determines which detections correspond to objects seen in frame $t-1$.

SORT (Simple Online and Realtime Tracking) provides the baseline approach, using Kalman filtering to predict object positions and the Hungarian algorithm for optimal detection-to-track assignment. However, SORT relies solely on motion information, causing frequent identity switches when objects pass close together.

DeepSORT addresses this by adding appearance information: a CNN extracts feature vectors capturing visual characteristics like clothing color or vehicle type. Combining motion predictions with appearance similarity significantly reduces identity switches in crowded scenes.

ByteTrack introduced an effective innovation: retaining low-confidence detections for secondary association. Low-confidence detections often occur when objects are partially occluded, yet they represent real objects. ByteTrack prevents premature track termination by attempting to associate these weak detections with existing tracks, achieving 80.3\% Multiple Object Tracking Accuracy (MOTA) and 77.3\% Identity F1 Score (IDF1) at 30 FPS.

Deep-OC-SORT (Deep Observation-Centric SORT) represents the current state-of-the-art method, ranking first on the MOT20 benchmark as of 2024. It combines the observation-centric paradigm where new detections drive the association process with deep appearance features. Deep-OC-SORT achieves 81.2\% IDF1 on MOT17, a 3.9\% improvement over ByteTrack, with only 1.8 milliseconds additional latency per frame. For forensic investigation tasks where accuracy matters more than real-time speed, this trade-off is highly favorable.

\subsection{Motion Prediction with Adaptive Kalman Filtering}

The Kalman filter estimates object state (position and velocity) by recursively updating predictions as new measurements arrive. The standard constant-velocity model assumes linear motion, which fails during speed changes, stops, or abrupt direction changes all common in urban traffic.

The Adaptive Kalman Filter learns motion model parameters from observed data, adjusting to object behavior patterns. Rather than assuming constant velocity, it incorporates information about accelerations, decelerations, and typical speed profiles. Empirical studies show the adaptive approach improves accuracy by 0.56\% HOTA and 0.81\% MOTA with negligible computational overhead. Because all cameras in the target deployment are stationary (mounted on fixed infrastructure), no camera motion compensation is needed, simplifying the system design.

\section{Background on Person and Vehicle Re-Identification}

Single-camera tracking produces trajectory segments (tracklets) within each camera's view. To create complete multi-camera trajectories, the system must identify when a person or vehicle seen in one camera matches an entity seen in another camera. This task re-identification (Re-ID)is challenging due to dramatic appearance changes across cameras.

\subsection{The Re-Identification Challenge}

The same person or vehicle appears different across cameras due to intra-class variation: different viewpoints (front vs. rear), lighting conditions (sunlight vs. artificial light), and camera angles change appearance. Inter-class similarity describes the opposite problem: different individuals may look similar (many white sedans, many people in dark clothing). Camera-specific biases compound the problem different sensors exhibit different color calibrations.

Standard benchmarks for evaluating Re-ID include Market-1501 (1,501 person identities across 6 cameras), DukeMTMC-reID (1,404 identities across 8 cameras), VeRi-776 (776 vehicles across 20 cameras), and MARS (video person Re-ID with 17,503 tracklets). Performance is measured using Rank-k accuracy (percentage of queries where the correct match is among the top $k$ results) and mean Average Precision (mAP).

\subsection{Deep Learning Approaches for Re-Identification}

Modern Re-ID systems employ deep neural networks to extract appearance features (embeddings) such that embeddings of the same identity are close together in feature space while embeddings of different identities are far apart.

Vision Transformers (ViTs) offer superior performance to CNNs by replacing convolutions with self-attention mechanisms. TransReID applies a Vision Transformer backbone to Re-ID with two key innovations: Side Information Embedding (SIE) encodes the camera index as learnable parameters, allowing the network to learn camera-specific appearance transformations; Jigsaw Patch Module (JPM) trains the network on shuffled image patches, encouraging learning of robust part-level features. TransReID achieves 95.2\% Rank-1 on Market-1501 (person Re-ID) and 86.1\% Rank-1 on VeRi-776 (vehicle Re-ID).

For video sequences processing multiple frames of a tracklet rather than single frames VID-Trans-ReID extends TransReID with temporal attention. Rather than processing each detected object frame-by-frame, it aggregates information across multiple frames of a tracklet, producing more robust embeddings. VID-Trans-ReID achieves 96.36\% Rank-1 accuracy on MARS (video person Re-ID), outperforming image-based methods by 1-2\%. For vehicle re-identification, applying temporal aggregation to vehicle tracklets provides similar robustness benefits, especially valuable for Egyptian CCTV that often operates at lower resolution with variable lighting.

\subsection{Feature Post-Processing and Re-Ranking}

After extracting deep features, post-processing techniques improve Re-ID accuracy without retraining. PCA (Principal Component Analysis) whitening decorrelates feature dimensions and normalizes their variances, improving cosine similarity discriminability by 5-7\%.

The k-Reciprocal Encoding method exploits query-gallery relationships. If query $q$ is among the $k$ nearest neighbors of gallery sample $g$, and $g$ is among the $k$ nearest neighbors of $q$, they are k-reciprocal neighbors, suggesting high match confidence. The Jaccard distance computed on these reciprocal sets refines similarity scores. Extensive evaluation shows k-reciprocal re-ranking provides +13.99\% mAP improvement on Market-1501, with optimal hyperparameters ($k_1=20$, $k_2=6$, $\lambda=0.3$) validated across multiple benchmarks.

\section{Comparative Study of Recent Work}

This section reviews state-of-the-art methods in multi-camera vehicle tracking, person re-identification, and system deployments from 2022 to 2025.

\subsection{Multi-Camera Vehicle Tracking}

The AI City Challenge, hosted at CVPR (Computer Vision and Pattern Recognition), is the premier competition for multi-camera tracking. The CityFlow dataset comprises 46 synchronized cameras across real-world intersections with 880 annotated vehicle identities, providing a rigorous benchmark.

Li et al., ranking second in the 2022 AI City Challenge with 0.8437 IDF1, proposed a comprehensive system incorporating domain knowledge. Their approach divides each camera's field of view into spatial zones and merges tracklets that exit one zone and re-enter from an adjacent zone within a short time window, repairing tracklets fragmented by occlusion. For multi-camera matching, they enforce directional constraints based on road topology: vehicles are grouped by movement direction and only compared with tracklets in adjacent cameras traveling the same direction, reducing false matches by 50-80\%.

The Camera Link Model (CLM) addresses manual setup burden by self-supervised learning. CLM identifies adjacent camera pairs by analyzing high-confidence tracklet matches, then learns zone pairings and transition time distributions automatically. This approach requires no manual camera topology annotation, making it adaptable to new deployments.

\subsection{Person Re-Identification Research}

Zhu et al. (2024) proposed camera-aware re-identification methods that explicitly handle camera-specific appearance biases. Their approach normalizes feature distributions per camera, embeds camera index information, and uses a camera-center triplet loss robust to label noise. This achieves state-of-the-art results on MSMT17 (15 cameras), demonstrating that camera-aware training improves cross-camera matching.

Multimodal fusion combining body and face features achieves robustness improvements. Studies report 89.1\% Rank-1 accuracy on Market-1501 and 79.8\% on DukeMTMC-reID by jointly using body appearance and facial features.

\subsection{Latest Tracking Benchmarks}

Current benchmarks demonstrate continued progress. Deep-OC-SORT ranks first on MOT20 with 63.9\% HOTA (Higher Order Tracking Accuracy) and second on MOT17 with 64.9\% HOTA as of 2024. This combination of observation-centric association and deep appearance features achieves the best identity consistency.

\subsection{Vector Similarity Search Infrastructure}

For forensic investigation, efficient similarity search over millions of Re-ID vectors is critical. FAISS (Facebook AI Similarity Search), developed by Meta AI Research, is the industry standard library. FAISS supports multiple indexing strategies: exhaustive brute-force search (slow but exact), Inverted File Index (IVF) with Product Quantization (PQ) (approximate but fast), and Hierarchical Navigable Small World (HNSW) graphs. FAISS provides GPU acceleration, achieving 8.5× speedup over CPU for billion-scale datasets. Meta uses FAISS in production for similarity search over billions of vectors, confirming its reliability and scalability.

\section{Implemented Approach and Technology Selection}

This section explains the specific technologies chosen for this project and the reasoning behind each selection.

\subsection{System Architecture Overview}

The implemented system processes video archives through a seven-stage pipeline. Stage~0 handles frame ingestion and preprocessing. Stage~1 performs per-camera object detection and single-camera tracking. Stage~2 extracts appearance embeddings and colour features from each tracklet. Stage~3 indexes all embeddings in a vector database alongside SQLite metadata. Stage~4 performs cross-camera association to assign globally consistent identities. Stage~5 evaluates tracking performance against ground truth. Stage~6 produces annotated visualisations and exports.

The architecture emphasises modularity: each stage produces standardised outputs (JSON tracklet files, NumPy embedding arrays, binary FAISS index, SQLite database) consumed by the next stage. A FastAPI backend and a Next.js frontend provide the operator-facing application layer on top of the pipeline.

\subsection{Object Detection: YOLO26m}

YOLO26m (Ultralytics) is selected for object detection. It is deployed using pre-trained COCO weights detecting the vehicle classes: car, bus, and truck (COCO class IDs 2, 5, and~7). The model runs at 1280-pixel input resolution to maximise detection of small vehicles in high-resolution traffic footage, with a confidence threshold of 0.25 and NMS IoU of 0.65 tuned for dense vehicle scenes. Inference uses half-precision (FP16) on GPU for efficiency. The SPPF architectural module in the backbone significantly improves detection of small and partially occluded objects common in urban intersections.

\subsection{Single-Camera Tracking: Deep-OC-SORT}

Deep-OC-SORT is selected for single-camera tracking because it achieves state-of-the-art accuracy: 81.2\% IDF1 on MOT17, representing a 3.9\% improvement over ByteTrack. For forensic investigation tasks where manual review of trajectories is necessary, higher identity consistency reduces operator effort.

The observation-centric paradigm, where new detections drive the association process, proves more robust to fragmented tracklets caused by occlusion common in crowded urban scenes (markets, transit hubs, congested intersections). The 1.8 millisecond increase in latency compared to ByteTrack is negligible for offline batch processing of recorded video.

The tracker is deployed via the BoxMOT library using OSNet-x0.25 appearance weights trained on MSMT17. A track buffer of 450 frames ($\approx$45 seconds at 10~FPS) is used to re-associate vehicles that stop at traffic lights without losing their track identity. Detection and tracking confidence thresholds are 0.25 and 0.05 respectively. After tracking, a TrackletBuilder post-processor applies linear bounding-box interpolation (maximum gap 30 frames) and within-camera intra-merge (maximum temporal gap 25 seconds, minimum tracklet area 800 pixels).



\subsection{Re-Identification: TransReID ViT-Base/16 CLIP}

TransReID with a Vision Transformer ViT-Base/16 backbone pre-trained on CLIP is selected for Re-ID. The CLIP pre-training provides rich semantic vehicle and person representations learned from image-text pairs at internet scale, giving the backbone a stronger initialisation than ImageNet-only supervised pre-training.

The model is fine-tuned on the CityFlowV2 dataset using Kaggle P100 GPU infrastructure. Input images are resized to 256$\times$256 pixels. The network produces 768-dimensional embeddings, which are subsequently compressed to 384 dimensions via PCA whitening to decorrelate feature dimensions and improve cosine similarity discriminability. All embeddings are L2-normalised before storage and retrieval.

Side Information Embedding (SIE) encodes the camera index as learnable parameters, allowing the network to learn camera-specific appearance transformations — directly addressing the colour calibration variance across different camera manufacturers. The Jigsaw Patch Module (JPM) trains on shuffled image patches, encouraging the network to learn robust part-level vehicle features rather than holistic appearance.

A secondary OSNet-x0.25 model (torchreid library) contributes a 10\% score-level fusion weight alongside the primary 90\% TransReID score, providing complementary lightweight features. The system achieves [PLACEHOLDER]\% mAP and [PLACEHOLDER]\% Rank-1 accuracy on the CityFlowV2 benchmark.

\subsection{Appearance Features: HSV Color Histograms}

HSV color histograms are used for fast coarse filtering. The 512-dimensional histogram (16 Hue × 8 Saturation × 8 Value) with center-region weighting is computed in parallel with deep feature extraction. During forensic search, color-based queries ("find all red vehicles") retrieve candidate tracklets in milliseconds, then deep features are applied to the narrow candidate set. This two-stage approach dramatically reduces search latency from minutes to seconds for 10-100 million tracklet archives.

\subsection{Feature Post-Processing: PCA Whitening and Spatial-Temporal Validation}

Two post-processing stages enhance Re-ID accuracy without retraining. PCA whitening reduces the 768-dimensional TransReID output to 384 dimensions, decorrelating feature dimensions and improving cosine similarity discriminability by 5--7\%. Spatial-temporal validation removes candidate cross-camera matches with implausible transition times based on known camera distances and typical vehicle speeds. K-reciprocal re-ranking was evaluated during development but found to provide no improvement on the CityFlowV2 benchmark and is disabled in the final system.

\subsection{Multi-Camera Association: Conflict-Free Connected Components}

For multi-camera association, conflict-free connected components clustering using NetworkX is implemented. This approach constructs a similarity graph where tracklets are nodes and pairs with weighted similarity above a threshold are connected by edges. Weighted similarity fuses appearance (TransReID cosine similarity, weight 0.70), HSV colour histogram similarity (weight 0.25), and spatial-temporal plausibility (weight 0.05). Top-K=20 mutual nearest neighbours are considered per tracklet query.

Connected components — maximal subgraphs where all nodes are reachable from each other — are extracted, each assigned a unique global track ID. The conflict-free constraint prevents a single camera's tracklet from appearing in two different components, avoiding false transitive identity chains. Final hyperparameters: similarity threshold 0.53, maximum component size 12 tracklets.

This method is unsupervised (no labelled association data required), computationally efficient ($O(V+E)$ time complexity), and interpretable: investigators can inspect the similarity graph and understand why two tracklets are assigned the same global identity.

\subsection{Vector Similarity Search: FAISS}

FAISS IndexFlatIP is used for exact cosine similarity search over the 384-dimensional L2-normalised embeddings. Exact search (IndexFlatIP) is preferred over approximate methods such as IVF-PQ because the system prioritises recall over memory efficiency: in forensic retrieval, missing a true match is more costly than increased search time. FAISS is developed by Meta AI Research and is the industry standard for high-performance embedding retrieval, used in production over billions of vectors.

\subsection{Web Application Stack}

A full-stack web application is implemented on top of the pipeline to provide an operator-facing interface. The backend is a FastAPI (Python) REST API serving 15 routers covering video upload, frame serving, detection results, tracklet data, pipeline execution control, timeline editing, cross-camera search, result export, and system health monitoring. Metadata is persisted in SQLite. The backend exposes a CORS-enabled API consumed by the frontend over HTTP.

The frontend is built with Next.js~14 and TypeScript, styled with Tailwind~CSS and the shadcn/ui component library (Radix~UI primitives). Zustand manages client-side application state. TanStack Query~v5 handles server data fetching and cache invalidation. Video.js provides in-browser video playback with quality level support. dnd-kit enables drag-and-drop timeline editing for tracklet review.

\section{Summary}

This chapter reviewed the foundational technologies for multi-camera tracking and re-identification. Section~3.1 established the principles of object detection using YOLO26m and single-camera tracking using Deep-OC-SORT. Section~3.2 explained how visual features are extracted and matched across cameras using TransReID ViT-Base/16 CLIP for re-identification. Section~3.3 reviewed recent research advances from 2022 to 2025.

Section~3.4 justified the technology selections for this project. The choice of YOLO26m for detection, Deep-OC-SORT via BoxMOT for tracking, TransReID ViT-Base/16 CLIP for re-identification, FAISS IndexFlatIP for vector search, conflict-free connected components via NetworkX for multi-camera association, and a FastAPI and Next.js~14 application layer represents a practical balance between accuracy, computational efficiency, and interpretability.